Search throughput uses the synchronous host-to-host boundary in
\cref{fig:3-boundary}, including query staging, transfers and preprocessing,
such as JHQ's rotation and CAGRA int8 query quantisation. All methods
upload through pinned (page-locked) host buffers; JHQ returns results to
ordinary pageable memory and the baselines to pinned memory. C++ drives
JHQ/RaBitQ and Python drives CAGRA/IVF-PQ. Handles and output buffers
are reused. These interface costs are
included and matter for small batches and near ties.

\par\medskip\noindent\begin{minipage}{\linewidth}
\centering\captionsetup{hypcap=false}
\resizebox{\linewidth}{!}{% Method-neutral search-call boundary; not a stage-time plot.
\begin{tikzpicture}[
 font=\small\linespread{1}\selectfont,
 box/.style={draw=ThesisNavy!65,rounded corners=2pt,align=center,
 text width=2.9cm,minimum height=1.7cm,inner sep=4pt},
 arr/.style={-{Latex[length=2mm]},draw=ThesisNavy!80,thick}]
\node[box] (host) at (0,0) {\mbox{Host FP32 queries}\\\mbox{Pinned staging}};
\node[box,fill=ThesisNavy!7] (upload) at (3.65,0)
 {\mbox{H2D transfer}\\Query\\\mbox{preprocessing}};
\node[box,fill=ThesisNavy!7] (gpu) at (7.3,0)
 {GPU search\\kernels\\\mbox{(method-specific)}};
\node[box] (result) at (10.95,0)
 {D2H results\\GPU\\\mbox{synchronisation}};
\draw[arr] (host) -- (upload);
\draw[arr] (upload) -- (gpu);
\draw[arr] (gpu) -- (result);
\draw[ThesisNavy,thick] (-1.6,1.15) -- (-1.6,1.45)
 -- (12.55,1.45) -- (12.55,1.15);
\node[fill=white,inner sep=3pt] at (5.475,1.45)
 {Search timing: host to host};
\node[box,text width=6.2cm,minimum height=1.5cm] (setup) at (1.75,-2.45)
 {Outside search timing\\Data loading; model training/loading\\Index construction; warm-up};
\node[box,text width=6.2cm,minimum height=1.5cm] (cal) at (9.2,-2.45)
 {Budget calibration (Chapter 6)\\Timed separately from search};
\draw[densely dashed,ThesisNavy!60] (-1.6,-1.1) -- (12.55,-1.1);
\end{tikzpicture}
}
\captionof{figure}[Host and device timing boundaries]{Timing starts with FP32 queries in host memory and ends with
results back on the host. It includes input staging, host-to-device (H2D)
transfer, query preprocessing, GPU search, device-to-host (D2H) transfer
and synchronisation. Construction, warm-up and budget calibration are
excluded; construction and calibration costs are reported separately.
GPU stage times come from separate profiling runs and need not sum to
this end-to-end interval.}
\label{fig:3-boundary}
\end{minipage}\par\medskip

The main comparison uses batches of 1,000, five measurement rounds and normally 24 calls
per round after warm-up. For batch size $b$, round $r$ has $m_r$ calls with
durations $t_{ri}$ in seconds. Throughput divides the batch size by the
mean call time within each round, then reports the median across rounds:
\begin{equation}
 \bar t_r=\frac1{m_r}\sum_{i=1}^{m_r}t_{ri},
 \qquad Q_r=\frac b{\bar t_r},
 \qquad Q=\operatorname{median}_r Q_r.
 \label{eq:3-throughput}
\end{equation}
The mean retains valid slow calls. Results measure synchronous batch
throughput, rather than concurrent-request tail latency. Rounds share an
index; separate cold builds assess training variation (\cref{sec:7-7}).
Warm-up and setup details are in \cref{sec:a-e1-protocol}.

The main comparison ran on an RTX 5090 (32,607\,MiB; 575\,W power limit), driver 595.71.05,
CUDA 13.0 and cuVS 26.08.01. OpenMP used 16 threads with 25 logical CPUs
available. JHQ's kernels target compute\_80 PTX and are compiled for the GPU by
the driver before timing \cite{cuda_platform}. GPU clocks and background host activity were not
controlled and may contribute to timing variation. Follow-up experiments
ran in later sessions. Their reports identify which methods were timed
together, whether trained models were reused and which arithmetic
precision was used
(\cref{sec:a-supplement,sec:4-2,sec:7-precision}).

Memory reporting distinguishes index allocation, whole-process search
allocation, including context and workspace, and construction peak.
JHQ's net figure subtracts a background allocation; CAGRA's figure measures
the index allocation, so their boundaries differ (\cref{sec:4-4,sec:7-7}).
Construction peaks are sampled every 20\,ms and may miss shorter
transients. Construction time covers train/add calls, excluding input
loading.
